% from ega3-1-fr.tex
\oldpage[0\textsubscript{III}]{56}
on note ensuite que l'on a
$\Gamma(X',\mathcal L'^\bullet)=\Gamma(X,\psi_*(\mathcal L'^\bullet))$ et
que l'homomorphisme composé
\[
\Gamma(X,\mathcal L^\bullet)\longrightarrow
\Gamma(X',\psi^*(\mathcal L^\bullet))
\xrightarrow{\Gamma(v)}\Gamma(X',\mathcal L'^\bullet)
\]
n'est autre que
\[
\label{0.12.1.7.2-fr}
\tag{12.1.7.2}
\Gamma(v^\flat):\Gamma(X,\mathcal L^\bullet)
\longrightarrow\Gamma(X,\psi_*(\mathcal L'^\bullet))
\]
(\hyperref[0.3.7.1-fr]{0\textsubscript{I}, 3.7.1}), et
(\hyperref[0.12.1.7.1-fr]{12.1.7.1}) s'obtient par passage à la cohomologie
dans (\hyperref[0.12.1.7.2-fr]{12.1.7.2}). D'autre part, les suites spectrales
du foncteur composé $\mathcal F'\leadsto\Gamma(X,\psi_*(\mathcal F'))$
s'obtiennent en considérant une résolution injective de Cartan-Eilenberg
$\mathcal M^{\bullet\bullet}=(\mathcal M^{ij})$ du complexe
$\psi_*(\mathcal L'^\bullet)$ dans la catégorie des faisceaux de groupes
abéliens sur $X$; les suites spectrales en question sont celles du bicomplexe
$\Gamma(X,\mathcal M^{\bullet\bullet})$ (qui sont birégulières puisque
$\mathcal M^{ij}=0$ pour $i<0$ ou $j<0$). Or, la première suite spectrale de
ce bicomplexe \emph{dégénère}, car les faisceaux
$\psi_*(\mathcal L'^i)$ sont flasques (G, II, 3.1.1), donc
$H_{\mathrm{II}}^q(\Gamma(X,\mathcal M^{i,\bullet}))
=H^q(\psi_*(\mathcal L'^i))=0$ pour $q>0$ (G, II, 4.4.3); on a donc des
edge-homomorphismes \emph{bijectifs} (\hyperref[0.11.1.6-fr]{11.1.6})
\[
\label{0.12.1.7.3-fr}
\tag{12.1.7.3}
'E_2^{i0}=H^i(H_{\mathrm{II}}^0(\Gamma(X,\mathcal M^{\bullet\bullet})))
\longrightarrow H^i(\Gamma(X,\mathcal M^{\bullet\bullet}))
\]
et on sait (\hyperref[0.11.3.4-fr]{11.3.4}) que cet homomorphisme provient,
par passage à la cohomologie, de l'augmentation
\[
\label{0.12.1.7.4-fr}
\tag{12.1.7.4}
\Gamma(X,\psi_*(\mathcal L'^\bullet))
\longrightarrow\Gamma(X,\mathcal M^{\bullet\bullet})
\]
laquelle provient elle-même de l'augmentation
$\eta:\psi_*(\mathcal L'^\bullet)\to\mathcal M^{\bullet\bullet}$. D'autre
part, pour la seconde suite spectrale, on a des edge-homomorphismes
\[
\label{0.12.1.7.5-fr}
\tag{12.1.7.5}
''E_2^{i0}=H^i(H_{\mathrm I}^0(\Gamma(X,\mathcal M^{\bullet\bullet})))
\longrightarrow H^i(\Gamma(X,\mathcal M^{\bullet\bullet}))
\]
provenant (\hyperref[0.11.3.4-fr]{11.3.4}) par passage à la cohomologie, de
l'homomorphisme de complexes
$Z_{\mathrm I}^0(\Gamma(X,\mathcal M^{\bullet\bullet}))
\to\Gamma(X,\mathcal M^{\bullet\bullet})$. Or, comme $\psi_*$ est exact à
gauche, la suite
\[
0\longrightarrow\psi_*(\mathcal F')\longrightarrow
\psi_*(\mathcal L'^0)\longrightarrow\psi_*(\mathcal L'^1)
\]
est exacte; par définition d'une résolution de Cartan-Eilenberg
(\hyperref[0.11.4.2-fr]{11.4.2}), on peut donc prendre
$B_{\mathrm I}^0(\mathcal M^{\bullet\bullet})=0$,
$Z_{\mathrm I}^0(\mathcal M^{\bullet\bullet})=\mathcal L^\bullet$; comme le
diagramme
\[
\xymatrix@C=50pt@R=24pt{
\mathcal L^0\ar[r]^{i^0}&\mathcal M^{00}\\
\psi_*(\mathcal F')\ar[u]^\varepsilon
\ar[r]_{\varepsilon'}\ar[ur]^{\varepsilon''}
&\psi_*(\mathcal L'^0)\ar[u]_{\eta^0}
}
\]
est commutatif, l'injection de complexes
$i:\mathcal L^\bullet\to\mathcal M^{\bullet\bullet}$ est compatible avec les
augmentations $\varepsilon$ et $\varepsilon''$. On a ainsi deux
homomorphismes de complexes de $\mathcal L^\bullet$ dans
$\mathcal M^{\bullet\bullet}$
\[
\xymatrix@C=52pt@R=22pt{
\mathcal L^\bullet\ar[r]^i\ar[dr]_{v^\flat}
&\mathcal M^{\bullet\bullet}\\
&\psi_*(\mathcal L'^\bullet)\ar[u]_\eta
}
\]
